% ============================================================
% Chapter 7 — Sensing, Logging, Replay, and Coding
% ============================================================
\chapter{Sensing, Logging, Replay, and Retrospective Coding}
\label{ch:sensing}

\roadmap{This chapter documents the platform's instrumentation: the
webcam facial pipelines and the exact model-to-signal attribution; the
eye-gaze and pupil channels; the two affective-state derivations; the
area-of-interest attention layer and its \allocscore{}; the logging
inventory and its privacy posture; the time-alignment anchor; and the
session-replay and retrospective-coding surfaces. It states each signal's
validity plainly---distinguishing measured, derived, and exploratory
channels---and closes with the known data-quality caveats that constrain
interpretation.}

\section{Overview and validity tiers}
\label{sec:sensing-overview}

The sensing layer captures, per learning session, five families of
signal: webcam facial analysis (emotions and action units), webcam eye
gaze, an exploratory pupil-size proxy, dense interaction telemetry, and
screen-state snapshots for replay. It is useful to state up front the
\emph{validity tier} of each, because the report's claims are calibrated
to them: the facial and interaction channels are measured signals from
established tools; the affective-state and AOI layers are \emph{derived}
interpretations with documented assumptions; and the pupil channel is an
\emph{exploratory proxy} of low validity. \Cref{fig:sensing-pipeline}
shows the facial pipeline, whose model attribution is specified next
because project documentation historically confused it.

\begin{figure}[htbp]
\centering
\begin{tikzpicture}[
  node distance=6mm and 10mm,
  box/.style={draw, rounded corners=2pt, align=center, minimum height=8mm,
              inner sep=4pt, font=\small},
  store/.style={box, fill=blockgreen},
  proc/.style={box, fill=blockblue},
  wrk/.style={box, fill=blockamber},
  arr/.style={-{Stealth[length=2.2mm]}, semithick}
]
  \node[proc] (cam) {Webcam segment\\\scriptsize 640$\times$480, 15\,fps, no audio};
  \node[store, right=of cam] (minio) {MinIO};
  \node[proc, right=of minio] (redis) {Redis queues};
  \node[wrk, above right=4mm and 12mm of redis] (of3) {OpenFace 3.0 worker};
  \node[wrk, below right=4mm and 12mm of redis] (pyf) {Py-Feat worker};
  \node[store, right=14mm of of3] (emo) {\dbtable{emotion\_frames}\\\scriptsize 8-class probs};
  \node[store, right=14mm of pyf] (au) {\dbtable{pyfeat\_au\_results}\\\scriptsize 18 AU intensities};
  \draw[arr] (cam) -- (minio);
  \draw[arr] (minio) -- (redis);
  \draw[arr] (redis) -- (of3);
  \draw[arr] (redis) -- (pyf);
  \draw[arr] (of3) -- node[above,font=\scriptsize]{5\,fps} (emo);
  \draw[arr] (pyf) -- node[below,font=\scriptsize]{1\,fps} (au);
\end{tikzpicture}
\caption[Facial sensing pipeline]{The facial sensing pipeline. Two
independent workers consume the same webcam segment: OpenFace~3.0 produces
eight-class emotion probabilities into \dbtable{emotion_frames} at 5\,fps;
Py-Feat produces eighteen action-unit intensities into
\dbtable{pyfeat_au_results} at 1\,fps.}
\label{fig:sensing-pipeline}
\end{figure}

\section{Facial analysis: definitive model attribution}
\label{sec:sensing-facial}

The attribution below was audited directly against the two workers'
inference code, and the report states it without hedging because a CSV
header in the export code mislabels one signal (\cref{sec:sensing-caveats}).

\begin{itemize}
  \item \textbf{Eight-class emotion probabilities come from OpenFace~3.0,
  not Py-Feat.} The OpenFace~3.0 worker runs genuine multitask inference
  (RetinaFace alignment \citep{deng2020retinaface} plus a multitask head) and writes per-frame
  AffectNet-ordered probabilities---neutral, happiness, sadness,
  surprise, fear, disgust, anger, contempt---together with the dominant
  class into \dbtable{emotion_frames} at 5\,fps.\footnote{\code{apps/openface3-worker/inference.py};
  \code{apps/openface3-worker/main.py}}
  \item \textbf{Eighteen action-unit intensities come from Py-Feat.} The
  Py-Feat worker runs its detector with the \code{xgb} AU model and writes
  eighteen AU intensities (AU01, 02, 04, 05, 06, 07, 09, 10, 12, 14, 15,
  17, 20, 23, 24, 25, 26, 28) into \dbtable{pyfeat_au_results} at
  1\,fps.\footnote{\code{apps/pyfeat-worker/processor.py};
  \code{apps/pyfeat-worker/db.py}} Py-Feat's own emotion model is loaded
  but never read or persisted, so Py-Feat contributes nothing to
  \dbtable{emotion_frames}.
  \item \textbf{Head pose is never computed.} \dbtable{emotion_frames} has
  yaw/pitch/roll columns, but the worker hard-codes them null and discards
  the model's gaze tensor. The report makes no head-pose
  claim.\footnote{\code{apps/openface3-worker/main.py:176};
  \code{inference.py:104}}
\end{itemize}

Both workers perform real inference and are inert only if deployed without
their model weights---an operational, not scientific, limitation.

\section{Eye gaze}
\label{sec:sensing-gaze}

Gaze is estimated in the browser by WebGazer
\citep{papoutsaki2016webgazer} using a weighted-ridge regression over
facial-landmark features, sampled at 10\,Hz and buffered for batched
upload.\footnote{\code{apps/web/src/lib/webgazer/useWebgazer.ts}} Each
\dbtable{WebgazerLog} row stores viewport coordinates, an optional
confidence, the page path, and a timestamp; confidence is carried through
untouched and later used as a quality filter in AOI scoring. Calibration
runs on new sessions and can be re-prompted after inactivity. The channel
is region-level: it is adequate to say which \emph{panel} a learner looks
at, not which \emph{word}, and all gaze-based claims respect that bound.

\section{Pupil size: an exploratory proxy}
\label{sec:sensing-pupil}

The pupil channel is presented as \emph{exploratory} because it is not
pupillometry. It borrows the WebGazer camera stream, samples at 2\,Hz, and
estimates a ``diameter'' by a dark-pixel-area heuristic over a fixed
central crop of the
frame:\footnote{\code{apps/web/src/lib/pupil-size/usePupilSize.ts}}
\begin{equation}
  d = 2\sqrt{\dfrac{n_{\text{dark}}}{\pi}},
\end{equation}
where $n_{\text{dark}}$ is the count of sub-threshold (dark) pixels in the
crop. Because the crop is a fixed rectangle rather than a localized iris,
$d$ conflates pupil, eyelash, brow, shadow, and nostril. The signal
derives from real webcam pixels and is not random, but its construct
validity as pupil diameter is low; the report treats any pupil-derived
quantity as an exploratory proxy and excludes it from primary claims.

\section{Affective-state derivation: two unwired pipelines}
\label{sec:sensing-affect}

The platform derives four learning-centered states---engagement, boredom,
confusion, frustration---from emotion probabilities in \emph{two}
independent places that are not unified, and the report describes both.

\subsection{Server-side windowed pipeline}
The persisted pipeline slides a window (default 30\,s, 10\,s stride,
minimum five faces per window) over \dbtable{emotion_frames},
mean-pools the eight probabilities across face-bearing frames, and applies
a configurable, versioned rule set to produce the four state scores and a
dominant state, stored in
\dbtable{AffectiveStateWindow}.\footnote{\code{apps/api/src/affective-mapping};
\code{packages/shared/src/affective-mapping.ts}} The default rules are
Wickens-informed weightings---for example, engagement as a weighted sum
favoring neutral, happiness, and surprise; boredom as a weighted sum over
neutral, sadness, and disgust attenuated by an arousal modifier;
frustration as a max over disgust/anger/(sadness,contempt); confusion as a
conjunction of anger, disgust, and (surprise,fear). Teachers may edit
these rules, and edits are version-tracked.

\subsection{Client-side per-frame classifier}
Independently, the replay tab and CSV exporter classify each emotion frame
with a \emph{different}, hard-coded threshold scheme---continuous
engagement and boredom scores and binary confusion/frustration flags---
whose weights differ from the server's (e.g., engagement
$0.4\,\text{happy}+0.4\,\text{neutral}+0.2\,\text{surprise}$ versus the
server's $0.45/0.35/0.20$).\footnote{\code{.../tabs/ReplayTab.tsx};
\code{.../lib/exportReplayCsv.ts}} The replay tab's threshold sliders
re-label \emph{this} client-side timeline only; they never recompute the
server windows. The report flags this divergence explicitly: any figure
must state which pipeline produced it, and the two are not to be pooled.

\section{Area-of-interest attention layer}
\label{sec:sensing-aoi}

The AOI layer answers a SEEV-style question (\cref{sec:bg-seev}): does the
learner look where the task expects? Interface panels carry a
\code{data-replay-region} label; on each replay snapshot their bounding
rectangles are captured in viewport pixels into
\dbtable{SessionReplaySnapshot.aois}.\footnote{\code{apps/web/src/lib/interaction-log/useSessionReplayRecorder.ts}}
The scoring kernel walks gaze samples against snapshots.\footnote{\code{apps/web/src/pages/teacher/student-logs/lib/aoiScoring.ts}}

A gaze sample is admitted only if it passes a validity mask: confidence at
least $0.5$, not during a tab-hidden interval, not within a
$750\,\text{ms}$ post-navigation settle, and not inside an idle epoch.
Admitted samples are bucketed to the smallest containing panel (the
lesson and PDF panels roll up to a single \code{lesson+pdf} bucket).
The session is segmented into epochs by inferred activity, of which three
are scored: reading, intervention-active, and chatbot-dialogue. Each
scored epoch carries an \emph{expected} allocation over the four
buckets---ordinal SEEV/expectancy--value weights supplied as
placeholders---normalized to sum to one.

Per epoch, the \emph{observed} allocation is the fraction of admitted
gaze samples in each bucket, and the alignment between observed and
expected is one minus half their total-variation distance:
\begin{equation}
  \text{alignment}_{e}
  = 1 - \tfrac{1}{2}\sum_{b \in \mathcal{B}}
    \bigl| o_{e,b} - x_{e,b} \bigr|,
  \qquad \mathcal{B} = \{\text{sidebar, lesson+pdf, chatbot, header}\},
  \label{eq:alignment}
\end{equation}
where $o_{e,b}$ and $x_{e,b}$ are the observed and expected shares for
bucket $b$ in epoch $e$; the score lies in $[0,1]$, with $1$ perfect
alignment, and is undefined for idle, navigation, empty, or zero-weight
epochs. The session-level \allocscore{} is the duration-weighted mean of
the defined epoch alignments:
\begin{equation}
  \text{allocation\_score}
  = \frac{\sum_{e} \text{alignment}_{e}\,\Delta t_{e}}
         {\sum_{e} \Delta t_{e}},
  \label{eq:allocscore}
\end{equation}
over epochs $e$ with a defined alignment and duration $\Delta t_{e}$.
Because the expected weights are placeholders rather than fitted SEEV
parameters, \allocscore{} is reported as a relative attention-alignment
index, not a validated SEEV measure---a scope caution carried into
\cref{ch:limitations}.

\section{Logging inventory and privacy posture}
\label{sec:sensing-logging}

Client logging mounts an interaction logger (eight raw streams) and a
replay recorder (snapshots).\footnote{\code{apps/web/src/components/LoggingProvider.tsx};
\code{apps/web/src/lib/interaction-log/useInteractionLogger.ts}}
\Cref{tab:logging} inventories the streams with their cadence. Buffered
streams flush every 30\,s (and on page hide); snapshots are taken every
3\,s plus on navigation/visibility events and flushed every 10\,s.
Structured domain events are recorded to \dbtable{ActivityLog} under a
43-value action enum (reproduced in \cref{app:activityenum}), fire-and-forget
so that instrumentation cannot break the learning flow.

\begin{table}[htbp]
\centering
\small
\caption{Raw interaction-logging streams and cadence.}
\label{tab:logging}
\begin{tabular}{@{}p{0.2\linewidth}p{0.45\linewidth}p{0.28\linewidth}@{}}
\toprule
Stream & Captured & Cadence \\
\midrule
\dbtable{cursor_logs} & pointer x/y, coarse target & 100\,ms throttle \\
\dbtable{click_logs} & x/y, selector, text ($\leq$50 chars) & per click \\
\dbtable{scroll_logs} & scroll position and percent & 200\,ms throttle \\
\dbtable{keystroke_logs} & \emph{counts}, pause, words/min & per field, on blur \\
\dbtable{clipboard_logs} & action and text \emph{length} & per event \\
\dbtable{visibility_logs} & visible/hidden, hidden duration & per event \\
\dbtable{viewport_logs} & size, orientation & init + resize \\
\dbtable{performance_logs} & load and API timings & per navigation \\
\dbtable{error_logs} & message, truncated stack & on error \\
\dbtable{session_replay_snapshots} & screenshot/DOM, geometry, AOIs & 3\,s + events \\
\bottomrule
\end{tabular}
\end{table}

The privacy posture is a design commitment verified in code: keystrokes
are stored as \emph{counts and typing speed, never key contents}; the
clipboard stores \emph{lengths, never contents}; password fields are
redacted in any DOM snapshot; element identifiers are coarse (tag, id, or
first class) and captured click text is capped; and webcam capture
records \emph{no audio}. What is deliberately not captured is as much a
result of the design as what is.

\subsection{Time alignment}
All streams are unified on one time base. At session start the client
records a \dbtable{session_sync_anchors} row tying wall-clock,
monotonic-clock, and (server-stamped) receive time together; the replay
and every export express time relative to this anchor's
wall-clock.\footnote{\code{useInteractionLogger.ts}; \code{logs.service.ts}}
Because facial, gaze, and pupil rows all carry wall-clock timestamps, the
modalities align without per-stream offset estimation (a
\dbtable{modality_offsets} table exists but is unused).

\section{Session replay}
\label{sec:sensing-replay}

The replay tab reconstructs a session as a scrubber-driven playback of
screen state, synchronized webcam video, and signal
overlays.\footnote{\code{apps/web/src/pages/teacher/student-logs/tabs/ReplayTab.tsx}}
To remain responsive on long sessions it loads in three stages: session
metadata with all interaction and biometric streams; then paginated
snapshot metadata (geometry and AOIs, no heavy content); then, on demand,
the full HTML and screenshot for the current snapshot, cached with
bounded size. Screen state is rendered in a sandboxed double-buffered
iframe with multi-pass scroll restoration, including a PDF-page fallback
for canvases still decoding. With DOM capture off---the study default---
the pixel screenshot is the primary view. Overlaid on the playback are the
gaze point, click pulses, the pupil proxy, the AOI rectangles, and a
timeline of the eight emotions, four learning states, and eighteen action
units.

\section{Retrospective coding}
\label{sec:sensing-coding}

Layered over replay is a qualitative coding
surface.\footnote{\code{apps/api/src/replay-annotations};
\code{.../components/AnnotationsPanel.tsx}} A researcher defines a private
code book (\dbtable{ReplayCode}: label and color) and attaches
time-anchored annotations (\dbtable{ReplayAnnotation}: a start, an optional
end, an optional code, and a note) expressed in the same anchor-relative
time as the replay. Annotations are additive---they never mutate captured
data---and per researcher, so several coders can annotate the same session
independently. This surface is the platform's provision for human-coded
ground truth; using it to obtain BROMP-style labels
\citep{ocumpaugh2015} is designed-but-future work (\cref{ch:future}), and
the report is careful to call the current validation \emph{convergent},
not criterion.

Deletion semantics are deliberate and matter for data retention: purging a
session cascades to its snapshots, activity log, recording segments, and
CV frames, and to any researcher annotations on it; but chatbot
conversation rows survive with a nulled session key. There is no scheduled
purge job---deletion is cascade-driven from session or user removal.

\section{Research exporters}
\label{sec:sensing-export}

Three export surfaces serve analysis. A wide-format replay CSV emits one
row per metric across a per-bin time grid (default 1\,s bins), with three
sampling regimes---nearest-sample for time-aligned signals (emotions,
action units, gaze, pupil, affective labels), carry-forward for continuous
state (scroll, viewport, AOI geometry), and bin-drop for sparse events
(clicks, keystrokes, utterances, EF detections)---and a catch-all row so
no activity action is silently lost. A full-session ZIP bundles a JSON
record, the webcam segments, and screenshot frames. A separate cohort
exporter in the studio server emits seven CSVs, including a split of
free-dialogue versus learning-strategy utterances, per-intervention
time-on-task, self-report rows, and EF detections, all honoring the
researcher's session-trim window. The column dictionary is given in
\cref{app:csvdict}.

\section{Known data-quality caveats}
\label{sec:sensing-caveats}

Five caveats constrain interpretation and are disclosed here so the
results can be read correctly.

\begin{enumerate}
  \item \textbf{Emotion source mislabeled in the CSV.} The replay CSV
  labels the emotion rows ``py-feat,'' but the data are OpenFace~3.0
  (\cref{sec:sensing-facial}). The values are correct; only the header
  string is wrong.\footnote{\code{exportReplayCsv.ts:392}}
  \item \textbf{Docked-layout scroll signal.} Before the scroll-host
  capture was added (mid-2026), window scroll was a dead signal in the
  docked course layout, so sessions recorded earlier cannot have their
  inner reading position reconstructed; the affected columns are null for
  those sessions. \todo{State the cutoff date and list affected session
  ranges once the deployment log is consulted.}
  \item \textbf{Biometric frame truncation.} A replay query once capped
  emotion and action-unit frames at 5{,}000 rows ($\approx$17 minutes at
  5\,fps); the cap was raised to 50{,}000 on 2026-05-21. Sessions longer
  than $\sim$17 minutes recorded before that date were silently truncated
  to their earliest frames. \todo{Identify affected sessions by date and
  duration.}
  \item \textbf{Hex-escape batch drops.} Malformed byte-escape fragments
  in browser-generated selectors can be rejected by the database's JSON
  parser; the ingest defensively drops the whole batch and returns success
  to avoid a stuck client buffer, so affected batches are lost silently
  (only a server warning records them) and downstream CSV bins have
  gaps.\footnote{\code{apps/api/src/logs/logs.service.ts}}
  \item \textbf{No automatic CV retry.} The Py-Feat job has no retry
  mechanism and no retry column; the OpenFace job's retry count increments
  only on manual re-runs. A failed job appears as missing biometric rows
  for that segment, with no automatic recovery; there is no observable
  automatic failure/retry rate to report.
\end{enumerate}

\medskip
\noindent These caveats are engineering realities, not fatal flaws, but
they bound what the data can support---which is exactly the subject of the
next chapter.
